\begin{afterword}
\summaryhead

The essay compares two ways of judging a claim: by who says it and by the argument given. A famous geologist is believed more than a troubled teenager, and a detailed physics argument more than a weak one. But the two behave differently. Once a listener has heard both speakers make equally good full arguments, the geologist's reputation adds little. Once the listener knows two arguers have equal credentials, the better argument still wins.

The author explains the difference with causal graphs. If night causes the sprinkler to run and the sprinkler makes the sidewalk slippery, then knowing the sprinkler is on makes the time of day irrelevant; the sprinkler ``screens off'' the night. Likewise truth produces arguments, and arguments produce expert belief, so knowing the argument screens off the expert. In practice, the author grants, some weight remains for authority, since experts may know unstated counterevidence or have hard-won intuitions. Algebra, he concludes, trumps authority.

\responsehead

The essay's conclusion is contained in its premises, and its own ending takes the premises back.

The thought experiments assume what they claim to find. The speakers give ``the best explanation they can possibly give,'' arguments that include ``all of the inferential steps'' and ``all of the support'' they used, to a listener ``knowledgeable enough to understand all the technical arguments.'' Under those conditions the speaker has nothing left to contribute by definition, and the essay says as much: ``A good technical argument is one that eliminates reliance on the personal authority of the speaker.'' The causal graph then encodes the same assumption as a picture. Truth leads to argument, argument leads to expert belief, and there is no other path. Screening off follows from the picture, as it would from any picture drawn that way.

The last four paragraphs supply the missing paths. Experts know counterevidence they did not mention. Their judgment of a step may rest on ``intuitions you can't duplicate without the same thirty years of experience.'' Each of these is an arrow from the truth to the expert that bypasses the argument, and with any of them the argument no longer screens off the expert. The essay calls this ``practice'' and keeps its title. On its own account the title holds for arguments reducible to pure mathematics, which are almost none of the arguments a reader receives from an expert.

Consider what the essay lets the reader do. It tells anyone who feels they have understood an argument that the arguer's credentials no longer matter. The listener who most needs to hear the opposite, the one who has not understood as much as they think, cannot tell from the inside that they are that listener. The essay's framing encourages the confidence: the low-authority speaker is a delinquent with psychotic episodes, the objector is a ``novice,'' and whatever the reader cannot follow is ``impressive noises.''

The author does not live by the rule. He asks the reader to take half the technical demonstration ``on my personal authority.'' He sends the reader to Pearl with advice about which book to read. His one weighing of authorities sets himself, at twenty-eight, a ``slightly'' lower probability than Jaynes. And the essay ends with a claim, an erratum in one of Jaynes's books, that the reader can accept only on the author's word.

In short: the title is true of a diagram drawn to make it true, the essay concedes it is false of real experts, and the author closes by claiming, on his own authority, to have corrected Jaynes.
\end{afterword}
